\documentclass{article}
\usepackage{hyperref}
\usepackage[a4paper]{geometry}
\usepackage{fancyhdr}
\pagestyle{fancy}
\lhead{Gelelektrophorese}
\rhead{März 2026}
\begin{document}
\section{Gelelektrophorese} 
Die Gelelektrophorese kann genutzt werden um unterschiedliche DNA-Abschnitte miteinander zu vergleichen.
 
\subsection{Ablauf}
Hier werden die DNA-Abschnitte welche mit der PCR vervielfältigt wurden in Gel gefüllt, an welches Gleichstrom angeschlossen wird. Die DNA ist negativ, bewegt sicht also zur Anode. Hier bewegen sich kleinere DNA-Abschnitte schneller, kommen weiter.
 
Nachdem sich die Abschnitte fertig bewegt haben, wird Farbstoff hinzugegeben, welche die Abschnitte markiert, am Gel aber nicht haftet.
 
So ist letztendlich ein der längen der DNA-Abschnitte abhängiges, identifizierbares Muster entstanden, welches mit anderen verglichen werden kann. 
 
\subsection{PCR}
Eine \textit{\textbf{p}olymerase \textbf{c}hain \textbf{r}eaction} wird genutzt um kurze DNA-Sequenzen schnell zu vervielfältigen. Dies funktioniert ähnlich wie die \hyperref[Replikation]{DNA-Replikation} in der Zelle. 
 
Die PCR wird in drei Phasen aufgeteilt. 
\begin{description} 
 \item[Denaturierung] Bei hohen Temperaturen, fast $100\,^\circ C$, brechen die Wasserstoff-Brückenbindungen der Doppelstränge der DNA; sie werden in zwei Einzelstränge aufgeteilt.
 \item[Primer-Hybridisierung] Bei ungefähr $60\,^\circ C$ kann der Primer, welche den Start den DNA-Synthese festlegen, effizient an komplementäre DNA-Sequenz binden.
 \item[Polymerisation] Bei $70\,^\circ C$ setzt die Taq-Polymerase, eine Polymerase welche an hohe Temperaturen angepasst ist, am Primer an und binder die komplementären Nukleotide.
\end{description} 
So wird in einem Zyklus jeder Doppelstrang einmal kopiert. 
\end{document}